% Generated draft: requires mathematical and source review.
\documentclass[11pt]{article}
\usepackage[margin=25mm]{geometry}
\usepackage{fontspec}
\setmainfont{Latin Modern Roman}
\usepackage{amsmath,amssymb,mathtools,unicode-math}
\setmathfont{Latin Modern Math}
\usepackage{enumitem,needspace,xurl,hyperref}
\setlength{\parindent}{0pt}
\setlength{\parskip}{5pt}
\setlength{\emergencystretch}{3em}
\newcommand{\N}{\mathbb{N}}
\newcommand{\Z}{\mathbb{Z}}
\newcommand{\Q}{\mathbb{Q}}
\newcommand{\R}{\mathbb{R}}
\newcommand{\C}{\mathbb{C}}
\newcommand{\F}{\mathbb{F}}
\newcommand{\A}{\mathbb{A}}
\newcommand{\PP}{\mathbb{P}}
\newcommand{\E}{\mathbb{E}}
\newcommand{\eps}{\varepsilon}
\newcommand{\dd}{\,\mathrm{d}}
\begin{document}
\section*{UnsolvedMath ID: 30001882}
\textbf{Source problem code:} catalogue entry
\subsection{Definitions and mathematical statement}
Let \(r,t\) be integers with \(r\geq t\geq1\), and let \(V\) be a finite set with \(|V|=n\). An \(r\)-uniform family on \(V\) is a collection \(\mathcal F\subseteq\binom{V}{r}\), where \(\binom{V}{r}=\{A\subseteq V:|A|=r\}\). The supplied record describes the extremal quantity schematically as
\[
\operatorname{UF}_t(n,r)=\max\bigl\{|\mathcal F|:\mathcal F\subseteq\tbinom{V}{r},\ \mathcal F\text{ is }t\text{-union-free}\bigr\}.
\]
The requested problem is to determine the asymptotic behavior of \(\operatorname{UF}_t(n,r)\) for \(r\geq t\geq1\).

The supplied record does not define \(t\)-union-free, so the displayed extremal expression is a transcription of the intended quantity rather than a complete mathematical definition. In particular, the admissible subfamilies whose unions are compared, their permitted sizes, and any distinctness or disjointness requirements are unspecified. These conditions must be recovered from the source before the problem has an unambiguous formulation. The asymptotic regime and the desired precision are also unspecified: taking \(n\to\infty\) with \(r,t\) fixed is a possible interpretation, but is not explicitly stated. Thus this draft records the extremal question with its unresolved qualifications.
\subsection{Short English statement}
For integers \(r\geq t\geq1\), determine how large a \(t\)-union-free family of \(r\)-element subsets of an \(n\)-element set can be asymptotically. The supplied record leaves both the definition of \(t\)-union-free and the precise asymptotic regime unresolved.
\subsection{Sources}
\begin{itemize}
\item Supplied local UnsolvedMath registry record, permanent ID \(30001882\), title \emph{Asymptotics of Uniform Union-Free Families}, source problem code \texttt{catalogue entry}, in the Oberwolfach Reports Open Problems collection. No individual proposer is supplied; the proposed year is recorded as 2011.
\item Supplied report citation: \emph{Discrete Geometry (2012)}, Oberwolfach Reports, DOI \texttt{10.4171/owr/2011/44}, pp.\ 2459--2548. Source URL and literature reference labelled ``original OWR report'': \url{https://doi.org/10.4171/owr/2011/44}. The report citation's 2012 label and the 2011 date information are retained as supplied; their relationship has not been verified.
\item Literature assessment dated exactly \texttt{2026-08-21}: the supplied curation labels the problem open and states that displayed bounds leave the requested asymptotic unresolved. This is supplied dated curation, not an independently verified literature assessment. The complete supplied record contains neither those bounds nor the report excerpt defining the family condition.
\item The registry labels the statement \texttt{corrected\_verified}; its accompanying verification note says only that TeX notation was enclosed in math delimiters without changing mathematical content. This label does not establish that the missing mathematical definition has been verified.
\end{itemize}
\end{document}
